\subsection{Capability-Safety Anti-correlation}\label{sec:tsr_car}

The displacement analysis in Section~\ref{sec:fake_safety} reveals a striking pattern at the aggregate level: models with higher task capability systematically report lower collision avoidance rates. Pi0.5 achieves a TSR of 74.7\,$\pm$\,1.5\% but a CAR of only 26.8\,$\pm$\,1.6\%, while Pi0-based methods (LoRA baseline, Full-FT, and A1-LoRA) report TSR $\leq$ 0.5\% alongside CAR values of 82--91.5\%. The anti-correlation is consistent across all three Pi0.5 seeds (TSR 73--76\%, CAR 25--28\%) and across Pi0 variants regardless of whether an SDF loss is applied. We observe that models capable of executing tasks also move the robot into configurations where collisions become possible, while models that fail to act avoid collisions by default.

Figure~\ref{fig:tsr_car_scatter} plots all evaluated configurations in TSR-CAR space. Two clusters emerge. Pi0-family methods occupy the upper-left region: high CAR, near-zero TSR. As established in Section~\ref{sec:fake_safety}, over 99\% of their safe episodes are fake safe, placing this cluster in what we term the \emph{fake safety zone}. Pi0.5, by contrast, sits in the lower-right region: high TSR, low CAR. Its collisions reflect genuine robot motion through the workspace. The two clusters are separated by a gap that no evaluated method occupies, a region that would require both meaningful task execution and reliable obstacle avoidance.

This spatial structure has a direct implication for metric interpretation. A model that reports high CAR but near-zero TSR provides no evidence of collision avoidance capability; its safety score is an artifact of inaction. CAR is informative only when conditioned on sufficient task capability. We note that the TSR-CAR plane offers a more complete picture than either metric alone: the upper-right quadrant (high TSR, high CAR) represents genuinely safe and capable behavior, while the upper-left quadrant is accessible through action collapse without any safety mechanism.

Early results from a Pi0.5 model augmented with A1-level SDF training (10K-step checkpoint) suggest that this gap is not inherent. This configuration achieves TSR of 47--49\% with CAR of 56--58\% and a frozen-action rate of 0\%, placing it in the previously unoccupied central region of Figure~\ref{fig:tsr_car_scatter}. We defer detailed analysis of this model to Section~\ref{sec:a1_pi05}, but note here that its position confirms the anti-correlation is a property of current models, not a fundamental constraint. Genuine collision avoidance requires a model that both moves and avoids, a combination that raw CAR cannot distinguish from one that simply does not move.
